\documentclass[full-course-notes.tex]{subfiles}

\begin{document}
\lecturemeta{7}
  {Transition Diagrams and Lexer Implementation Strategies: Ad Hoc vs.\ Table-Driven Scanning}
  {CLO 1, CLO 3}
  {\S3.4 (Recognition of Tokens: transition diagrams for \texttt{relop}, identifiers/keywords,
   and whitespace; sketching a lexical analyzer directly from transition diagrams)}
  {[ET] Engineering a Compiler, \S2.2--2.3 (hand-written scanners); [AP] Modern Compiler
   Implementation, \S2.5}
\lecshort{Transition Diagrams; Ad Hoc vs.\ Table-Driven Lexers}
\setcounter{section}{0}
\makelecturetitle

\tikzset{
  st/.style={draw, circle, minimum size=0.85cm, font=\small, fill=nuLight, draw=nuNavy},
  fin/.style={st, double, double distance=2pt, fill=dbGreenBg, draw=dbGreen},
  trans/.style={thick},
}

\begin{coreidea}[Lecture Roadmap]
\begin{enumerate}
  \item \textbf{Transition diagrams}: a lighter-weight, one-pattern-at-a-time cousin of the
  DFAs built in Lecture~\ref{lec:6} -- exactly the tool DB \S3.4 uses to go from ``here is a
  regular expression'' to ``here is code,'' worked through on \texttt{relop}, identifiers/keywords,
  and whitespace.
  \item \textbf{Retraction}: what it means for a state to need to ``look one character too far
  ahead,'' and how that connects directly back to Lecture~\ref{lec:6}'s double-buffering
  scheme.
  \item \textbf{The ad hoc approach}: translating a set of transition diagrams straight into
  hand-written scanning code, one routine per token category, the way many real compilers'
  lexers are actually built.
  \item \textbf{Ad hoc vs.\ table-driven, head to head}: the two genuinely different ways to
  \emph{implement} a lexer -- hand-coded diagrams (today) vs.\ the generic, generated DFA
  simulation Lecture~\ref{lec:6} already built in full -- and when a real engineer reaches for
  each one.
\end{enumerate}
Lectures~\ref{lec:5}--\ref{lec:6} gave you the \emph{formal} machinery (NFA, subset
construction, minimization) that a lexer-generator tool runs automatically. This lecture asks a
different, very practical question: what if you are writing the scanner \emph{by hand}, without
a generator at all? DB introduces transition diagrams for exactly this reason, and -- not
coincidentally -- this is also where Project Phase~1 and Assignment~1 begin.
\end{coreidea}

\section{From Formal Automata to a Concrete Scanner}
\label{sec:l7-motivation}

\dbtag DB actually introduces transition diagrams (\S3.4) \emph{before} it formalizes NFAs and
DFAs -- historically, hand-drawn diagrams like these are how lexers were built for years before
subset construction and minimization were automated into generator tools. Having already done
Lectures~\ref{lec:5}--\ref{lec:6} in full formal generality, we are in a good position to see
transition diagrams for what they really are: \textbf{a deliberately simplified, by-hand DFA},
drawn separately for \emph{one} token category at a time, skipping the machinery
(\Sref{sec:l6-lexgen-combining}'s tagged-union combination, full subset construction) that a
generator needs to merge \emph{many} categories into a single automaton automatically.

\begin{coreidea}[Why Bother, Given Lecture~\ref{lec:6} Already Solved This?]
Lecture~\ref{lec:6} answered ``how does a lexer-generator tool build a scanner automatically.''
This lecture answers a different, equally real question: ``what does a programmer actually
\emph{write}, by hand, when no generator tool is used at all?'' Both questions matter --
production compilers are built both ways, sometimes even the \emph{same} compiler mixing both
strategies for different parts of its lexer (\Sref{sec:l7-hybrid-reality} returns to this).
Understanding the hand-written route also makes the generated route feel less like magic: every
diagram in this lecture is exactly the kind of small DFA fragment that subset construction
(\Sref{sec:l6-subset-construction}) would have produced for you, had you chosen to generate
instead of hand-write it.
\end{coreidea}

\section{Transition Diagrams for Individual Token Patterns}
\label{sec:l7-diagrams}

\begin{defbox}{Transition Diagram (DB \S3.4 Sense)}
A \textbf{transition diagram} here is a DFA-like picture built for \emph{one} token category at
a time: circles for states, labeled arrows for the character(s) that trigger each transition,
and one or more accepting states. It differs from the general DFAs of Lecture~\ref{lec:6} in
one specific way: some accepting states are marked with a \textbf{$*$}, meaning ``this state
accepts, but only after the scanner has already read one character \emph{too far} to know that
-- that extra character must be pushed back before the token is considered complete.'' This
extra bookkeeping step is called \textbf{retraction}, covered in \Sref{sec:l7-retraction}
below.
\end{defbox}

\subsection{Worked Example: Identifiers and Keywords}
\label{sec:l7-diagram-id}

\dbtag The simplest and most-reused diagram: one letter, then any number of letters or digits
-- exactly Lecture~\ref{lec:4}'s \texttt{id} regular definition
(\Sref{sec:l4-worked-example-identifiers}), just drawn as a picture instead of written as a
formula.

\begin{center}
\begin{tikzpicture}[node distance=2.6cm, >=Latex]
  \node[st] (s0) {$0$};
  \node[st, right=of s0] (s8) {$8$};
  \node[fin, right=of s8] (s9) {$9^{*}$};
  \draw[trans, ->] (-1.3,0) -- (s0);
  \draw[trans, ->] (s0) -- node[above, font=\small]{letter} (s8);
  \draw[trans, ->] (s8) to[out=135, in=225, looseness=6] node[left, font=\small]{letter, digit} (s8);
  \draw[trans, ->] (s8) -- node[above, font=\small, align=center]{other\\(retract)} (s9);
\end{tikzpicture}
\end{center}

\begin{examplebox}{Reading Diagram: Identifier/Keyword}
State~8 self-loops on every letter or digit -- exactly maximal munch, one character at a time.
The moment a character arrives that is \emph{neither} a letter nor a digit (a space, an
operator, end of file), that character cannot be part of the identifier, so state~9 accepts --
but state~9 is marked $*$, because the scanner already consumed that one extra, non-identifier
character to \emph{discover} the identifier was finished. \Sref{sec:l7-retraction} explains
exactly what happens to that character next. Once the lexeme is known, the reserved-word trick
from \Sref{sec:l3-beyond-db} decides \texttt{ID} vs.\ a specific keyword token via a symbol-table
lookup -- the diagram itself cannot tell \texttt{while} apart from \texttt{score}; only the
lookup afterward can.
\end{examplebox}

\subsection{Worked Example: \texttt{relop} (DB's Classic Diagram)}
\label{sec:l7-diagram-relop}

\dbtag DB's single most-cited transition-diagram example: the six relational operators
\texttt{<}, \texttt{<=}, \texttt{<>}, \texttt{=}, \texttt{>}, \texttt{>=}. It is worth tracing
by hand once, since it is the clearest possible illustration of \emph{why} retraction is
needed at all -- some branches consume exactly one extra character to confirm a longer match,
others don't.

\begin{center}
\begin{tikzpicture}[node distance=2.55cm, >=Latex]
  \node[st] (s0) {$0$};
  \node[st, above right=1.5cm and 2.3cm of s0] (sA) {$A$};
  \node[fin, right=of s0, yshift=0cm] (sB) {$B$};
  \node[st, below right=1.5cm and 2.3cm of s0] (sC) {$C$};

  \node[fin, right=2.55cm of sA] (sD) {$D$};
  \node[fin, below=1.1cm of sD] (sE) {$E$};
  \node[fin, below=1.1cm of sE] (sF) {$F^{*}$};

  \node[fin, right=2.55cm of sC] (sG) {$G$};
  \node[fin, below=1.1cm of sG] (sH) {$H^{*}$};

  \draw[trans, ->] (-1.6,0) -- (s0);
  \draw[trans, ->] (s0) -- node[above, sloped, font=\small]{$<$} (sA);
  \draw[trans, ->] (s0) -- node[above, font=\small]{$=$} (sB);
  \draw[trans, ->] (s0) -- node[below, sloped, font=\small]{$>$} (sC);

  \draw[trans, ->] (sA) -- node[above, font=\small]{$=$} (sD);
  \draw[trans, ->] (sA) -- node[right, font=\small]{$>$} (sE);
  \draw[trans, ->] (sA) -- node[right, font=\small, align=center]{other\\(retract)} (sF);

  \draw[trans, ->] (sC) -- node[above, font=\small]{$=$} (sG);
  \draw[trans, ->] (sC) -- node[right, font=\small, align=center]{other\\(retract)} (sH);
\end{tikzpicture}
\end{center}

\begin{center}
\renewcommand{\arraystretch}{1.3}
\small
\begin{tabularx}{\textwidth}{@{}l l X@{}}
\toprule
\textbf{State} & \textbf{Token, Attribute} & \textbf{Retract Needed?} \\
\midrule
$B$ & \texttt{RELOP}, EQ (\texttt{=}) & No -- \texttt{=} alone is never a prefix of a longer
\texttt{relop} in this language, so the very next character read is already known to belong to
whatever comes \emph{after} this token. \\
$D$ & \texttt{RELOP}, LE (\texttt{<=}) & No -- both characters of the lexeme were consumed and
both were needed. \\
$E$ & \texttt{RELOP}, NE (\texttt{<>}) & No -- same reasoning as $D$. \\
$F^{*}$ & \texttt{RELOP}, LT (\texttt{<}) & \textbf{Yes} -- reaching $F$ means the scanner read
\texttt{<} \emph{and then} some other character (to rule out \texttt{<=} and \texttt{<>|}) that
turned out not to belong to this token at all. \\
$G$ & \texttt{RELOP}, GE (\texttt{>=}) & No. \\
$H^{*}$ & \texttt{RELOP}, GT (\texttt{>}) & \textbf{Yes} -- same reasoning as $F$. \\
\bottomrule
\end{tabularx}
\end{center}

\begin{examfocus}
Notice \emph{exactly} which states need the asterisk, and why: only the two states ($F$, $H$)
reached by taking the ``other'' branch out of a two-character lookahead need retraction. States
reached by consuming a character that was itself part of a \emph{necessary}, meaningful
transition ($B$, $D$, $E$, $G$) never need it. A very common exam mistake is marking every
accepting state with a $*$ out of caution -- retraction is needed \emph{only} when a character
was read purely to \emph{rule out} a longer match and then turned out not to belong to the
token at all.
\end{examfocus}

\subsection{Worked Example: Whitespace}
\label{sec:l7-diagram-whitespace}

\dbtag One more diagram worth seeing explicitly, because it is the one diagram that produces
\emph{no token at all} -- exactly \Sref{sec:l3-tokens-patterns-lexemes}'s ``housekeeping''
point that whitespace and comments are consumed and discarded, not emitted as tokens.

\begin{center}
\begin{tikzpicture}[node distance=2.6cm, >=Latex]
  \node[st] (w0) {$0$};
  \node[fin] (w1) at ($(w0) + (2.6,0)$) {$1^{*}$};
  \draw[trans, ->] (-1.3,0) -- (w0);
  \draw[trans, ->] (w0) to[out=135, in=225, looseness=6] node[left, font=\small, align=center]{blank,\\tab,\\newline} (w0);
  \draw[trans, ->] (w0) -- node[above, font=\small, align=center]{other\\(retract)} (w1);
\end{tikzpicture}
\end{center}

\begin{examplebox}{Why Whitespace Still Needs Retraction}
State~0 self-loops through any run of spaces, tabs, or newlines. The instant a non-whitespace
character appears, that is exactly the retract situation again: the scanner has read one
character belonging to the \emph{next} real token, purely to discover whitespace has ended.
State $1^{*}$ ``accepts,'' but instead of emitting a token, the scanner simply \textbf{restarts}
matching from state~0 of whichever diagram applies to the retracted character -- whitespace is
the one pattern in a real lexer specification whose successful match produces zero tokens, not
one.
\end{examplebox}

\section{Retraction and the Buffer, Revisited}
\label{sec:l7-retraction}

\dbtag Every $*$ in the diagrams above means the same concrete thing:
\Sref{sec:l6-lexgen-maximal-munch}'s maximal-munch loop already described this exact situation
generically (``keep advancing... roll back to the last remembered accepting position''), and
\Sref{sec:l6-lexgen-buffering}'s double-buffering scheme is precisely what makes that rollback
cheap. Nothing new is being invented here -- transition diagrams just make the \emph{single
most common form} of that rollback (exactly one character of overshoot) visible and explicit,
state by state, instead of leaving it implicit inside a generic simulation loop.

\begin{coreidea}[Retraction, Concretely]
Reaching a $*$ state means: move the \texttt{forward} pointer (\Sref{sec:l6-lexgen-buffering})
back by exactly one position before recording where the lexeme ends. Because
\Sref{sec:l6-lexgen-buffering}'s buffers keep recently-read characters in memory, this is a
pointer decrement, not a re-read from disk -- exactly why double buffering was worth building
in the first place. In hand-written, ad hoc code (\Sref{sec:l7-adhoc} below), this is often
implemented even more directly: use a \textbf{non-consuming lookahead} (a \texttt{peekChar()}
that reads the next character \emph{without} advancing the input pointer) instead of an
advance-then-retract pair. The two are behaviorally identical -- ``peek, then only advance if
the character turns out to belong to this token'' never overshoots in the first place, so
there is nothing left to retract.
\end{coreidea}

\section{From Diagrams to Code: The Ad Hoc Approach}
\label{sec:l7-adhoc}

\dbtag With a transition diagram in hand for every token category, \textbf{ad hoc} lexer
construction is exactly what the name suggests: translate each diagram directly into ordinary
code, by hand, one routine (or one block of a larger routine) per category, dispatched by
inspecting the first character of whatever remains in the input.

\begin{defbox}{Ad Hoc Lexer}
An \textbf{ad hoc lexer} is a hand-written scanner in which each transition diagram becomes an
explicit block of code -- typically an \texttt{if}/\texttt{switch} dispatch on the first
character, followed by a small loop or a short chain of character tests that walks that one
diagram's states directly, with no shared, generic ``simulate a table'' loop underneath. The
programmer, not a generator tool, is responsible for combining the diagrams (deciding dispatch
order, handling the reserved-word trick, wiring in retraction) correctly.
\end{defbox}

\begin{examplebox}{A Small Ad Hoc Scanner, Sketched}
The three diagrams from \Sref{sec:l7-diagrams} translate almost line for line:
\begin{verbatim}
Token getNextToken() {
    skipWhitespace();          // Sec. l7-diagram-whitespace, inlined
    c = peekChar();            // non-consuming lookahead (Sec. l7-retraction)

    if (isLetter(c)) {         // Sec. l7-diagram-id
        lexemeBegin = forward;
        do { advance(); c = peekChar(); }
        while (isLetter(c) || isDigit(c));
        lexeme = getLexeme(lexemeBegin, forward);
        if (isKeyword(lexeme))
            return makeToken(keywordTokenFor(lexeme));
        return makeToken(ID, installID(lexeme));
    }

    if (c == '<') {            // Sec. l7-diagram-relop, the '<' branch
        advance();
        if (peekChar() == '=') { advance(); return makeToken(RELOP, LE); }
        if (peekChar() == '>') { advance(); return makeToken(RELOP, NE); }
        return makeToken(RELOP, LT);   // F*/H*: peek already IS the retraction
    }
    if (c == '=') { advance(); return makeToken(RELOP, EQ); }
    if (c == '>') {
        advance();
        if (peekChar() == '=') { advance(); return makeToken(RELOP, GE); }
        return makeToken(RELOP, GT);
    }

    // ... one block per remaining category (NUM, LITERAL, ...) ...

    reportLexicalError(c);
}
\end{verbatim}
Notice the $F^{*}$/$H^{*}$ states from \Sref{sec:l7-diagram-relop} never actually call an
explicit ``retract'' function here at all -- using \texttt{peekChar()} throughout means the
input pointer only ever moves forward past a character once that character is confirmed to
belong to the current token, so there is nothing to undo. This is the standard real-world way
ad hoc scanners sidestep explicit retraction bookkeeping, exactly as
\Sref{sec:l7-retraction}'s coreidea box described.
\end{examplebox}

\begin{examfocus}
Be able to translate a small transition diagram (given to you on an exam) directly into
pseudocode of this shape: one \texttt{if}/\texttt{switch} branch per outgoing edge from the
start state, a loop for any self-loop, and a clear point where each accepting state either
returns a token or (for whitespace) simply restarts the dispatch. This diagram-to-code
translation is exactly what Assignment~1 and Project Phase~1 ask you to do by hand.
\end{examfocus}

\section{The Table-Driven Approach, Revisited}
\label{sec:l7-tabledriven-recap}

\dbtag Lecture~\ref{lec:6} already built this approach in full, so this section is a short
recap, not new material: combine \emph{every} token pattern's NFA into one, via the tagged-union
construction (\Sref{sec:l6-lexgen-combining}); run subset construction
(\Sref{sec:l6-subset-construction}) and minimization (\Sref{sec:l6-minimization}) once, offline;
encode the result as a transition table; and drive that table with \emph{one generic loop} --
\Sref{sec:l6-lexgen-maximal-munch}'s maximal-munch simulation -- that has no idea, and does not
need to know, how many token categories it is handling. A lexer-generator tool like Flex
automates this entire pipeline (\Sref{sec:l6-lexgen-pipeline}) from a rules file you write, and
never once produces or reasons about a hand-drawn diagram like today's.

\begin{examfocus}
Keep straight which lecture built \emph{which} half: Lecture~\ref{lec:6} built the
\textbf{table-driven} half (one combined automaton, one generic simulation loop, machine-built
via subset construction). \emph{This} lecture built the \textbf{ad hoc} half (many small,
separately hand-drawn diagrams, each translated into its own bespoke code). Both are complete,
correct ways to end up with a working lexer -- they differ in \emph{how} the lexer is
constructed, not in what it ultimately does.
\end{examfocus}

\section{Ad Hoc vs.\ Table-Driven: A Head-to-Head Comparison}
\label{sec:l7-comparison}

\begin{center}
\renewcommand{\arraystretch}{1.45}
\begin{tabularx}{\textwidth}{@{}l X X@{}}
\toprule
 & \textbf{Ad Hoc (this lecture)} & \textbf{Table-Driven (Lecture~\ref{lec:6})} \\
\midrule
\textbf{How it's built} & By hand: draw one diagram per token category, translate each into its
own block of code & Automatically: combine all patterns' NFAs, run subset construction and
minimization once \\
\textbf{Combining many patterns} & The programmer's responsibility -- dispatch order,
reserved-word handling, and retraction must all be gotten right by hand for \emph{every}
diagram & Handled uniformly by the tagged-union construction
(\Sref{sec:l6-lexgen-combining}) and rule-order tie-breaking
(\Sref{sec:l6-lexgen-ties}) -- correct by construction \\
\textbf{Adding a new token} & Write (and correctly integrate) a new block of code, ordered
correctly relative to every existing block & Add one line to the rules file and regenerate --
the tool re-derives the whole combined automaton \\
\textbf{Performance} & Can be \emph{very} fast -- a human can special-case the hottest patterns
(e.g.\ single-character operators) with no generic overhead at all & Also fast ($O(1)$ per
character via table lookup, \Sref{sec:l6-nfa-vs-dfa-tradeoffs}), but pays for whatever the
combined table's size turns out to be \\
\textbf{Correctness guarantees} & None automatic -- a forgotten retraction, a keyword checked
after (not before) the generic \texttt{ID} branch, or a missed edge case is a real, easy-to-make
bug & Strong -- subset construction and minimization are provably correct algorithms; if the
rules file is right, the generated scanner is right \\
\textbf{Debuggability} & Ordinary source code -- step through it in any debugger, exactly like
any other hand-written function & A generic loop over an opaque table -- debugging usually means
inspecting \emph{table contents} and current state number, one level more abstract \\
\textbf{Best fit} & Small, fixed, performance-critical token sets where hand-tuning pays off, or
where no generator tool is available/desired & Larger or evolving token sets, and any situation
where the specification (the rules file) should be the single source of truth \\
\bottomrule
\end{tabularx}
\end{center}

\begin{enrichment}[Extra Depth: The Hybrid Reality]
\label{sec:l7-hybrid-reality}
Just as Lecture~\ref{lec:1} showed that ``compiled vs.\ interpreted'' is rarely a pure either/or
choice in practice (most real systems are hybrids), the same is true here. CPython's own
tokenizer and V8's (Chrome/Node.js JavaScript engine) lexer are both, in large part,
\textbf{hand-written} C/C++ code -- not output from Flex or a similar tool -- precisely because
a hand-tuned fast path for the handful of hottest patterns (identifiers, common operators,
string-literal boundaries) can outperform a fully generic table-driven loop, while still calling
out to smaller, more mechanically-generated pieces (e.g.\ Unicode category tables) for the parts
where hand-tuning would not pay off. GCC's own lexer has, at different points in its history,
used both strategies. The lesson: ``ad hoc vs.\ table-driven'' is a real engineering choice with
real trade-offs (as the table above lays out), not a solved question with one universally
correct answer -- exactly why this course teaches both.
\end{enrichment}

\begin{coreidea}[Where This Fits: Project Phase~1 and Assignment~1]
This lecture is deliberately timed to sit right before you start writing code of your own.
\textbf{Assignment~1} (regular expressions and finite automata) and \textbf{Project Phase~1}
(the lexer stage of your semester project, first due around Lecture~7 in the master course
outline) both draw directly on today's material: you will draw transition diagrams for your own
language's tokens, and then choose -- or be told to use -- one of the two implementation
strategies above to turn those diagrams into a running scanner, before Lecture~\ref{lec:6}'s
Opt.~1 hands-on session shows the alternative, generator-based route (Flex) side by side with
what you just built by hand.
\end{coreidea}

\section{Self-Check Questions}
\label{sec:l7-selfcheck}

\begin{enrichment}[Practice -- Not Graded]
\begin{enumerate}
  \item Draw a transition diagram, in the style of \Sref{sec:l7-diagram-relop}, for a
  hypothetical \texttt{arrow} token family: \texttt{-}, \texttt{->}, \texttt{--} (decrement).
  Mark every state that needs retraction, and justify each marking.
  \item For DB's \texttt{relop} diagram (\Sref{sec:l7-diagram-relop}), trace the input
  \texttt{<>=} character by character. Which token is emitted first, and what happens to the
  remaining \texttt{=}?
  \item Explain, without saying ``because DB says so,'' why state $B$ (matching bare \texttt{=})
  in \Sref{sec:l7-diagram-relop} never needs retraction, while states $F^{*}$ and $H^{*}$ always
  do.
  \item Translate the whitespace diagram (\Sref{sec:l7-diagram-whitespace}) into pseudocode in
  the same style as \Sref{sec:l7-adhoc}'s worked example, and explain why it never calls
  \texttt{makeToken(...)}.
  \item A classmate's ad hoc lexer checks the generic identifier pattern \emph{before} checking
  for specific keywords like \texttt{while}. What token does their lexer incorrectly emit for
  the lexeme \texttt{while}, and how does \Sref{sec:l3-beyond-db}'s reserved-word trick fix
  this, concretely, in code?
  \item Using \Sref{sec:l7-comparison}'s comparison table, argue \emph{for} using an ad hoc
  lexer for a small teaching language with 15 fixed token types, and separately argue
  \emph{against} using ad hoc for a language whose token set is still actively evolving (like an
  in-development scripting language). What changes between the two scenarios?
  \item Explain concretely how using a non-consuming \texttt{peekChar()} (\Sref{sec:l7-retraction})
  avoids ever needing an explicit ``move the pointer back one'' operation, even though every
  $*$-marked state in this lecture's diagrams conceptually represents exactly that situation.
\end{enumerate}
\end{enrichment}

\begin{takeaways}
\begin{itemize}
  \item A \textbf{transition diagram} is a small, by-hand DFA for \emph{one} token category;
  states marked $*$ are exactly the states Lecture~\ref{lec:6}'s maximal-munch loop reaches by
  reading one character past the true end of the token.
  \item \textbf{Retraction} means giving that one overshot character back, so scanning can
  resume from the right place -- cheap in practice because of Lecture~\ref{lec:6}'s
  double-buffered input, and often sidestepped entirely in hand-written code by using a
  non-consuming \texttt{peekChar()} lookahead instead.
  \item DB's classic \texttt{relop} diagram is the cleanest worked illustration of \emph{when}
  retraction is and isn't needed: only branches that read a character purely to \emph{rule out}
  a longer match, and then discover it doesn't belong, need it.
  \item The \textbf{ad hoc} approach translates each diagram directly into hand-written,
  bespoke code -- fast and fully debuggable, but every bit of correctness (dispatch order,
  keyword handling, retraction) is the programmer's own responsibility.
  \item The \textbf{table-driven} approach (Lecture~\ref{lec:6}, in full) combines every pattern
  into one automaton mechanically and drives it with one generic, provably-correct simulation
  loop -- easier to extend and to trust, at the cost of being one level more abstract to debug.
  \item Real production lexers (CPython, V8, historically GCC) often \textbf{mix both
  strategies}, hand-tuning the hottest patterns while generating or table-driving the rest --
  exactly the same "no single universally correct answer" lesson Lecture~\ref{lec:1} already
  drew for compiled vs.\ interpreted execution.
  \item This lecture is the direct on-ramp to \textbf{Assignment~1} and \textbf{Project
  Phase~1}: drawing your own transition diagrams and choosing an implementation strategy for
  your own language's tokens.
\end{itemize}
\end{takeaways}

\section*{References for This Lecture}
\begin{thebibliography}{99}
\bibitem[DB]{db} Aho, A.~V., Lam, M.~S., Sethi, R., and Ullman, J.~D. \textit{Compilers:
Principles, Techniques, and Tools}. 2nd~ed. Pearson, 2006.
\bibitem[ET]{et} Cooper, K.~D. and Torczon, L. \textit{Engineering a Compiler}. 2nd~ed. Morgan
Kaufmann, 2011.
\bibitem[AP]{ap} Appel, A.~W. \textit{Modern Compiler Implementation in C/Java/ML}. Cambridge
University Press, 2004.
\end{thebibliography}

\end{document}
